% Generated draft: requires mathematical and source review.
\documentclass[11pt]{article}
\usepackage[margin=25mm]{geometry}
\usepackage{fontspec}
\setmainfont{Latin Modern Roman}
\usepackage{amsmath,amssymb,mathtools,unicode-math}
\setmathfont{Latin Modern Math}
\usepackage{enumitem,needspace,xurl,hyperref}
\setlength{\parindent}{0pt}
\setlength{\parskip}{5pt}
\setlength{\emergencystretch}{3em}
\newcommand{\N}{\mathbb{N}}
\newcommand{\Z}{\mathbb{Z}}
\newcommand{\Q}{\mathbb{Q}}
\newcommand{\R}{\mathbb{R}}
\newcommand{\C}{\mathbb{C}}
\newcommand{\F}{\mathbb{F}}
\newcommand{\A}{\mathbb{A}}
\newcommand{\PP}{\mathbb{P}}
\newcommand{\E}{\mathbb{E}}
\newcommand{\eps}{\varepsilon}
\newcommand{\dd}{\,\mathrm{d}}
\begin{document}
\section*{UnsolvedMath ID: 800001}
\textbf{Source problem code:} catalogue entry
\subsection{Definitions and mathematical statement}
A semigroup is a nonempty set with an associative binary operation. An extension of a semigroup \(S\) is a semigroup \(T\) equipped with an injective homomorphism \(\iota:S\hookrightarrow T\). An equation with constants from \(S\) is an equality between nonempty words in variables and named elements of \(S\). A solution in an extension assigns elements of \(T\) to the variables, interprets each constant \(c\in S\) as \(\iota(c)\), and makes every equation hold simultaneously.

\textbf{Source problem (S. Kublanovsky and M. Sapir).} Describe the types of equations or systems of equations for which an algorithm can decide, for every choice of the constants in the original semigroup, whether the equation or system has a solution in some finite extension of that semigroup.

The following expresses the decision target conditionally, without resolving missing source conventions. For a finite system template with variables \(x_1,\ldots,x_n\), constant slots \(c_1,\ldots,c_r\), and equations
\[
\Sigma(\mathbf x;\mathbf c)=\{u_j(\mathbf x,\mathbf c)=v_j(\mathbf x,\mathbf c):1\leq j\leq m\},
\]
let \(\mathcal E_{\mathrm{fin}}(S)\) denote the extensions admitted by the intended, but unspecified, meaning of ``finite extension.'' For each \(\mathbf c\in S^r\), the required decision is whether
\[
\exists (T,\iota)\in\mathcal E_{\mathrm{fin}}(S)\;\exists\mathbf t\in T^n\;\bigwedge_{j=1}^{m}
 u_j(\mathbf t,\iota(\mathbf c))=v_j(\mathbf t,\iota(\mathbf c)).
\]
Here \(\iota(\mathbf c)\) means applying \(\iota\) to each constant. The requested classification concerns the equation or system types for which this decision can be made by an algorithm that halts on every admissible input.

\textbf{Unresolved qualifications.} The source does not define ``finite extension,'' specify the original semigroup or its effective presentation, or determine whether the algorithm is uniform over original semigroups, system templates, or both. It also does not explicitly restrict systems to finite systems; the displayed finite-template formulation is therefore a conditional clarification, not an additional source hypothesis. If ``finite extension'' means that \(T\) itself is finite, injectivity forces \(S\) to be finite. These choices must be settled before a unique algorithmic classification problem is specified.
\subsection{Short English statement}
Classify the equation and system types for which one can decide, for every assignment of constants from the original semigroup, whether a solution exists in some finite extension. The meaning of finite extension, the effective input model, and the required algorithmic uniformity remain unspecified in the supplied source.
\subsection{Sources}
\begin{itemize}
\item Local UnsolvedMath registry, permanent ID \(800001\), exact source code \texttt{catalogue entry}, title \emph{The Lyapin's Notebook --- Problem 1}; supplied upstream revision \texttt{d371ee917ac9647fbb06cd6fc529deec6f3dda82}. The registry labels the problem \texttt{open}; this is a supplied status, not an independently verified assessment.
\item Bershadsky, Kublanovsky, and Mashevitzky, \emph{The Lyapin's Notebook: Unsolved Problems in Semigroup Theory} (2026), Problem 1, \url{https://arxiv.org/abs/2604.04763}. The supplied statement attributes this problem specifically to S. Kublanovsky and M. Sapir. The registry's \texttt{proposed\_by} field instead lists Bershadsky, Kublanovsky, and Mashevitzky, with proposed year 2026; these registry fields do not establish the original proposal date or replace the statement's attribution.
\item Supplied provenance: extraction method \texttt{arxiv\_latex}, access date 2026-10-16. This access information is supplied provenance; the reference and its contents have not been independently checked for this draft.
\end{itemize}
\end{document}
